\section{Neo4j Services}

This section describes the Neo4j service layer responsible for managing the graph
representation of cars and implementing the recommendation system in Car-Rev.
The graph database is used for structural relationships and similarity-based
computations, while MongoDB stores the main document data.

\subsection{Graph Model Overview}

The graph database contains the following main node types:

\begin{itemize}
    \item \textbf{Car}
    \item \textbf{Brand}
    \item \textbf{Model}
    \item \textbf{BodyType}
    \item \textbf{DriveWheels}
    \item \textbf{Transmission}
    \item \textbf{FuelType}
    \item \textbf{Year}
    \item \textbf{PriceBin}
    \item \textbf{DispBin}
    \item \textbf{User}
\end{itemize}

The main relationships are:

\begin{itemize}
    \item \texttt{ (:Car)-[:OF_BRAND]->(:Brand)}
    \item \texttt{(:Car)-[:OF_MODEL]->(:Model)}
    \item \texttt{(:Car)-[:HAS_BODY_TYPE]->(:BodyType)}
    \item \texttt{(:Car)-[:HAS_DRIVE]->(:DriveWheels)}
    \item \texttt{(:Car)-[:HAS_TRANSMISSION]->(:Transmission)}
    \item \texttt{(:Car)-[:USES_FUEL]->(:FuelType)}
    \item \texttt{(:Car)-[:OF_YEAR]->(:Year)}
    \item \texttt{(:Car)-[:HAS_PRICE_BIN]->(:PriceBin)}
    \item \texttt{(:Car)-[:HAS_DISPLACEMENT_BIN]->(:DispBin)}
    \item \texttt{(:User)-[:HAS_VISITED]->(:Car)}
\end{itemize}

The bin nodes (PriceBin and DispBin) allow grouping cars into ranges in order to
improve similarity computation.

\subsection{Neo4jCarInsertService}

This service is responsible for creating the graph projection of a car when a new
vehicle is inserted into the system.

\subsubsection*{Responsibilities}

\begin{itemize}
    \item Create a \texttt{Car} node identified by \texttt{mongo\_id}.
    \item Normalize string attributes (lowercase standardization).
    \item Merge or create structural dimension nodes (Brand, Model, etc.).
    \item Attach bin nodes based on numeric ranges.
\end{itemize}

\subsubsection*{Design Rationale}

Instead of storing all attributes directly inside the Car node only,
categorical features are externalized into dedicated nodes. This allows:

\begin{itemize}
    \item Efficient similarity matching via shared neighbors.
    \item Reduced duplication of categorical values.
    \item Future graph-based clustering.
\end{itemize}

Price and displacement bins are matched using range-based queries,
ensuring each car belongs to at most one bin per category.

\subsection{Neo4jCarDeleteService}

This service removes a car projection from the graph database.

\subsubsection*{Responsibilities}

\begin{itemize}
    \item Match the \texttt{Car} node by \texttt{mongo\_id}.
    \item Remove the node using \texttt{DETACH DELETE}.
    \item Return the number of deleted nodes.
\end{itemize}

\subsubsection*{Consistency Considerations}

The graph projection is deleted independently from MongoDB.
The operation assumes that MongoDB deletion has already occurred or will occur
in a coordinated workflow.

Since Neo4j does not manage cascading consistency with MongoDB,
eventual consistency is accepted between the two systems.

\subsection{Neo4jUpdateService}

This service updates the graph projection of a car when its attributes change.

\subsubsection*{Responsibilities}

\begin{itemize}
    \item Match the existing Car node by \texttt{mongo\_id}.
    \item Update scalar properties.
    \item Remove outdated relationships.
    \item Recompute and reattach dimension and bin relationships.
\end{itemize}

\subsubsection*{Update Strategy}

The update follows a two-phase logic:

\begin{enumerate}
    \item Update node properties.
    \item Delete and recreate structural relationships.
\end{enumerate}

This ensures the graph remains aligned with the most recent
car specifications.

\subsubsection*{Trade-offs}

The deletion and recreation approach simplifies logic but
temporarily removes relationships during execution.
Given the transactional nature of Neo4j write operations,
this does not produce visible intermediate inconsistencies.

\subsection{Neo4jRecommendationService}

This service implements the content-based recommendation engine
using graph traversal and structural similarity.

\subsubsection*{Recommendation Logic}

For a given user:

\begin{enumerate}
    \item Retrieve visited cars.
    \item Use each visited car as a seed.
    \item Find candidate cars sharing structural nodes:
    \begin{itemize}
        \item same BodyType
        \item same DriveWheels
        \item same Transmission
        \item same FuelType
        \item same PriceBin
        \item same DispBin
    \end{itemize}
    \item Exclude already visited cars.
    \item Exclude cars of the same brand.
    \item Apply numeric tolerance constraints:
    \begin{itemize}
        \item price within ±20\%
        \item displacement within ±0.4L
    \end{itemize}
    \item Rank by:
    \begin{itemize}
        \item number of shared features (descending)
        \item price difference (ascending)
        \item displacement difference (ascending)
    \end{itemize}
    \item Return top-k results.
\end{enumerate}

\subsubsection*{Design Rationale}

The recommendation strategy is:

\begin{itemize}
    \item Content-based.
    \item Deterministic.
    \item Fully explainable.
\end{itemize}

Unlike collaborative filtering approaches,
this method does not rely on global user similarity
but on structural proximity in the feature graph.

\subsection{Architectural Role of Neo4j}

Neo4j is used exclusively for:

\begin{itemize}
    \item Feature-based similarity.
    \item User-visit tracking.
    \item Efficient traversal over categorical features.
\end{itemize}

All heavy document storage remains in MongoDB,
while Redis handles real-time counters and caching.

This separation ensures:

\begin{itemize}
    \item Scalability of read-heavy operations.
    \item Reduced duplication.
    \item Clear domain separation between storage models.
\end{itemize}
